\section{Anchor Dossier Model}

\subsection{Purpose}

The project now needs case files rather than broader rankings. This section turns the strongest candidate units into anchor dossiers: each dossier gathers occurrence rows, context, minimal alternations, component reuse, and allowed lexical search classes. The goal is controlled partial decipherment.

This remains pre-phonetic. The output proposes reading classes such as ``emblem-linked name/title candidate'' or ``administrative label candidate'', not sound values.

\subsection{Method}

The script \texttt{scripts/anchor-dossier-model.py} combines:

\begin{itemize}
  \item semantic reconstruction candidates;
  \item lexical reading gates;
  \item structural reconstructions and slot fillers;
  \item phonetic minimal-pair tests;
  \item abstract sign variables.
\end{itemize}

The model keeps all strong onomastic anchors and strong-context cross-frame anchors. It then asks whether each anchor is a stable unit, a compound, or a sequence of reusable morpheme-like components.

\subsection{Corpus-Wide Result}

\begin{table}[htbp]
\centering
\begin{tabular}{lr}
\toprule
\textbf{Metric} & \textbf{Value} \\
\midrule
Anchor dossiers & 11 \\
Strong onomastic dossiers & 7 \\
Cross-frame dossiers & 4 \\
Anchor occurrence rows & 59 \\
Direct minimal tests & 49 \\
Component minimal tests & 277 \\
Anchor components & 15 \\
Root-like components & 2 \\
Prefix-like components & 3 \\
Suffix-like components & 4 \\
\bottomrule
\end{tabular}
\caption{Anchor dossier summary.}
\label{tab:anchor-dossier-summary}
\end{table}

\subsection{Highest-Priority Dossiers}

\begin{itemize}
  \item \texttt{240-482}: institutional/admin label candidate; Harappa-heavy; mostly tablet context.
  \item \texttt{176}: emblem-linked name/title candidate; appears as \texttt{176} and in \texttt{240-176}.
  \item \texttt{235-222}: emblem-linked name/title candidate; bull context and Chanhu-daro concentration.
  \item \texttt{032-220}: emblem-linked name/title candidate; bull context across Mohenjo-daro and Harappa.
  \item \texttt{590-390}: emblem-linked name/title candidate; steatite seal context across three sites.
  \item \texttt{240-176}: controlled lexical candidate; may contain the same root-like \texttt{176} with \texttt{240} as a qualifier.
  \item \texttt{235-840-032}: controlled lexical candidate; shares \texttt{235} and \texttt{032} with other anchors.
\end{itemize}

\subsection{Morpheme-Like Component Probe}

The most important new result is component reuse inside the strongest anchors:

\begin{itemize}
  \item \texttt{176} appears as a standalone prime-entity filler and as the final component in \texttt{240-176}. This is the cleanest root-like entity component so far.
  \item \texttt{032} appears as a standalone cross-frame unit and inside \texttt{032-220} and \texttt{235-840-032}. This makes it root-like or subtype-marking, not merely decorative.
  \item \texttt{240} appears initially in \texttt{240-482}, \texttt{240-176}, and \texttt{240-740-090}. It is currently the strongest prefix/title/classifier-like component in the anchor set.
  \item \texttt{235} appears initially in \texttt{235-222} and \texttt{235-840-032}; it remains formula-marker or title-prefix-like.
  \item \texttt{220}, \texttt{222}, \texttt{390}, \texttt{482}, and \texttt{840} are lexical-root or final-qualifier candidates that need image-grounded validation.
\end{itemize}

\subsection{Working Partial-Reading Classes}

The strongest controlled reading classes are now:

\begin{center}
\texttt{240 + ENTITY} \quad and \quad \texttt{235 + ENTITY}
\end{center}

This does not mean \texttt{240} or \texttt{235} has a phonetic value. It means they behave like reusable entity qualifiers in the prime-entity zone. If future image review confirms that \texttt{176} and \texttt{240-176} share a semantic context, then \texttt{240} becomes a serious prefix/title/classifier candidate and \texttt{176} becomes a serious root-like candidate.

\subsection{Interpretation}

This is the first stage where the project can make a morpheme-like statement:

\begin{center}
\texttt{QUALIFIER / TITLE / CLASSIFIER + ENTITY STEM}
\end{center}

For example, \texttt{240-176} is no longer just a two-sign filler. It is a testable compound: \texttt{240} may qualify or classify the root-like \texttt{176}. The same logic applies to \texttt{235-222} and \texttt{235-840-032}.

This is a meaningful step toward language, because language identification will require reusable segments that behave predictably across compounds. We now have a small set of such segments, but they remain abstract until image context and minimal-pair constraints validate them.

\subsection{Outputs}

\begin{itemize}
  \item \texttt{outputs/anchor\_dossier\_summary.csv}
  \item \texttt{outputs/anchor\_dossiers.csv}
  \item \texttt{outputs/anchor\_occurrence\_evidence.csv}
  \item \texttt{outputs/anchor\_minimal\_evidence.csv}
  \item \texttt{outputs/anchor\_component\_roles.csv}
  \item \texttt{outputs/anchor\_component\_edges.csv}
  \item \texttt{outputs/anchor\_reading\_hypotheses.csv}
  \item \texttt{outputs/anchor\_dossier\_model.tex}
\end{itemize}
